% !TeX root = ../main_en.tex
\subsection{Conclusions}

This work has been developed in a reproducible and traceable environment to compare classic and neural imitation control for quadcopter trajectory tracking. The contribution consists of a complete pipeline: from a fully modular 6DOF simulator, to the creation of test scenarios, specialized PD tuning, MLP, GRU, and LSTM network training, and their closed-loop execution, up to the generation of comparative metrics.

The initial hypothesis is partially fulfilled, with nuances. Neural architectures trained to predict the desired force of the outer loop are viable on trajectory families known during training, and they can successfully replace the set of PD controllers in several cases and complete their objective. However, no global superiority of the neural networks is observed. Under out-of-distribution conditions, failures emerge, alongside notable differences between architectures and a clear dependency on the geometry and dynamic demand level of the reference trajectory.

The first conclusion addresses the specialization of classic control. Some transferred PDs perform better than the specialized PD for the specific family, while others lose practically all performance when departing from their tuning trajectory. This confirms that classic transfer is uneven, and therefore this set of specialized PDs represents a necessary baseline for evaluating neural imitation. It also shows that the PD tuning process could be refined in future versions, particularly in families where the specialized controller is not the best on its own trajectories or behaves better in transfer than in native mode.

The second conclusion concerns the comparison between supervised fidelity and closed-loop tracking. GRU and LSTM approximate the specialized PD's force better in supervised evaluation, but they do not systematically dominate position RMSE. The MLP, despite not using explicit temporal memory, obtains very competitive results in several ID families and in mixed OOD. Therefore, while supervised fidelity is necessary, it is not sufficient to predict the quality of dynamic tracking.

The third conclusion centers on recurrent memory, as it does not yield automatic improvement. The GRU offers a solid result in several simulations, whereas the LSTM exhibits severe failures on geometrically demanding trajectories that are faster than those encountered during training. This suggests that history can provide value in certain states, but can also introduce sensitivity to poorly represented states, larger delays, and accumulated deviations. Consequently, the choice of architecture should be analyzed in conjunction with the trajectory type and its specific characteristics.

The fourth conclusion concerns safety in simulations, which belongs to the control system as a whole, not just the neural network. The classic inner loop, mixer, actuator saturation, and force clipping condition the final behavior. This design decision was appropriate because it allowed studying the neural network's behavior while maintaining controllability over the rest of the system, although it required more thorough analysis before attributing observations to the neural component.

Overall, the work meets its objective of comparing the replacement of a set of PDs with a neural implementation of desired force, all within a traceable, reproducible, and coherently scoped experimental environment. The results show that neural imitation is promising under conditions known during training and in certain OOD variations, although generalization cannot be taken for granted.

\subsection{Scope and Limitations}

The physical scope of the work is defined by the 6DOF simulator described in Section~\ref{subsec:modelo-simulador}. The vehicle is represented as a rigid body; position and velocity are expressed in the East--North--Up (ENU) world frame, while angular velocity, moments, and rotor geometry use the Forward--Right--Down (FRD) body frame. Attitude is propagated using quaternions, and dynamics use Newton--Euler equations integrated with fixed-step RK4. The model includes gravity, simplified linear drag, constant wind, Gaussian noise for position and velocity observations, a quadratic rotor law, delay, first-order dynamics, and actuator saturation.

These choices represent modeling assumptions. Constant mass, geometry, inertia, and rotor parameters are assumed. Structural flexibility, ground effect, rotor-to-rotor interference, rotor flapping, non-linear aerodynamics, turbulence, battery discharge, and thrust variation with voltage are not modeled. Wind is treated as a constant disturbance, and linear drag as a simplified dissipative term. Observations incorporate synthetic noise in position and velocity; the expected variability of an IMU, GNSS receiver, or specific delays of real sensors is not included.

The scope of the experiments is also bounded. The in-distribution dataset contains 150 episodes distributed among \texttt{hold}, \texttt{circle}, \texttt{lissajous}, and \texttt{waypoint}, with 105 training episodes, 22 validation episodes, and 23 test episodes. Within this environment, ID families, PD transfer between families, and a separate OOD series of ten scenarios with compositions, variations, helices, and lemniscatas are evaluated. The summary in Table~\ref{tab:cobertura-ejecucion} collects 184 runs in \texttt{test} and 80 OOD runs, sufficient to test the hypothesis, though not to exhaustively cover all geometries, speed profiles, vertical actions, disturbances, or maneuvers that would occur in real flight.

The controller comparison is limited to fixed-parameter specialized PDs and three neural network architectures for external force prediction: MLP, GRU, and LSTM. The networks are trained by supervised behavior cloning of the desired force calculated by specialized PDs, preserving the classic inner attitude loop, mixer, saturations, and force and inclination limits. Therefore, a full neural policy or a controller trained to minimize position error in closed loop via reinforcement learning is not evaluated. The neural force prediction implementation can inherit biases from the PDs and from the limited coverage of the training set. Furthermore, the PD used as a reference in OOD scenarios is an approximation based on geometric similarity, not a controller optimized for those geometries.

The parametric study has remained limited. The PDs are tuned using a reproducible procedure and then their parameters are fixed, but no global search in the space of possible gains or exhaustive multi-objective optimization was run. Similarly, MLP, GRU, and LSTM share the input configuration, normalization method, number of hidden units, batch size, learning rate, early stopping criteria, and seed to isolate the effect of the architecture. Sensitivity tests were conducted on the number of hidden units, window size, and select seeds, but not a full sweep of hyperparameters, regularization, dataset size, or checkpoint selection criteria. No sweep over the physical parameters of the vehicle, delay and saturation levels, noise, or wind was performed either.

The statistical validity of the results is approached with caution. The runs are deterministic for a given scenario and seed, favoring reproducibility, but confidence intervals or estimates of variability across seeds are not provided.

Finally, attitude limits, persistent saturation, and force and inclination limits define the region of validity for the runs. These constraints are necessary for simulation analysis but condition interpretation: part of the behavior originates from the complete control architecture rather than the neural network alone.

\subsection{Future Work}

A first line of work consists of adapting the simulator to a specific vehicle. This requires identifying mass, inertia, rotor geometry, actuator limits, thrust curve, delays, saturations, and sensor parameters from tests or precise technical documentation. Once system identification is completed, the results observed in simulations would bear a more direct relationship to the capability of a real quadcopter.

The second line is to increase the physical fidelity of the model. The current simulator includes a bounded set of phenomena, sufficient to conduct this comparison, but more advanced models of aerodynamics, ground effect, turbulence, rotor-to-rotor interference, battery discharge effects, thrust variation with voltage, and thermal, electrical, and flexible structural constraints can be incorporated. These extensions would allow studying whether controllers maintain stable behavior when actuation response varies unpredictably during a simulation or when effects currently estimated in a simplified manner appear.

It is also necessary to expand and refine the set of trajectories. In this work, few representative families have been used, but a flight-oriented implementation should increase the number of families, their geometric diversity, speeds, vertical profiles, and transitions between different trajectories. This refinement should be accompanied by a review of termination, thrust, tilt, and saturation limits to reflect the actual demands of a given vehicle type.

Another highly relevant line is the systematic study of neural architectures. In this work, MLP, GRU, and LSTM are compared under a series of common configurations to isolate the effect of the architecture. In a subsequent phase, the objective would be to find the best configuration for trajectory tracking, exploring width, depth, window length, normalization, input selection, regularization, seeds, training set size, and checkpoint selection criteria. It is feasible that each architecture requires a different size and window, as well as possessing different strengths, an aspect hinted at in the discussed results.

Within that line, the inference method for recurrent networks should also be studied. Currently, a window of 20 samples is reconstructed and processed in each control cycle without directly reusing the hidden states calculated in the previous cycle. An alternative would be to preserve these states between steps, both to reduce recalculation and to approach a more stateful inference. However, this modification would alter the relationship of the network with its operation, since the output would depend on a state propagated between internal cycles rather than just a window of samples calculated at each step.

Likewise, it is advisable to conduct a correlation study to identify which data subsets contribute to generalization capacity. It is not enough to add more trajectories; rather, one must analyze which families, parameters, and disturbances truly enrich the training distribution. This line of research would allow distinguishing whether OOD failures stem from geometric limitations, a lack of high-acceleration states, overly restrictive actuation limits, or an architecture incapable of representing the relationship between observation, reference, and desired force in those types of trajectories.

The natural subsequent step is to transfer the approach to a real vehicle. To do so, simply loading the network trained in simulation is insufficient; portability to a lower-level programming language is required, as well as guaranteeing that the real quadcopter provides the necessary state, reference, and actuation variables. It would also be necessary to align control, telemetry, and clock frequencies between the simulator and the vehicle, define a progressive testing procedure, and establish safe operating limits.

Finally, a closing line of work consists of using the simulator as a pre-flight validation platform. This would allow testing new trajectories, detecting regimes where the controller demands excessive forces, adjusting protection limits, and bounding the risk for real trials. In this way, the developed work would serve as a baseline to transition from a simulation for comparisons to an experimental evaluation on hardware.
